% ============================================================================================
% Cost-repair in a real RLVR training loop. Numbers via % AUTO-GENERATED by scripts/make_rlvr_macros.py — do not edit by hand.
% [PENDING] = experiment not yet complete; rerun the generator to fill.
\newcommand{\dialFiftyAccEnd}{60.0}
\newcommand{\dialFiftyAccStart}{53.9}
\newcommand{\dialFiftyCalls}{19769}
\newcommand{\dialFiftyGameEnd}{19.2}
\newcommand{\dialFiftyGameStart}{21.8}
\newcommand{\dialFiftyHoBase}{66.0}
\newcommand{\dialFiftyHoFinal}{67.0}
\newcommand{\dialFiftyStrongCalls}{19769}
\newcommand{\dialSeventyFiveAccEnd}{61.2}
\newcommand{\dialSeventyFiveAccStart}{55.0}
\newcommand{\dialSeventyFiveCalls}{29736}
\newcommand{\dialSeventyFiveGameEnd}{18.9}
\newcommand{\dialSeventyFiveGameStart}{21.3}
\newcommand{\dialSeventyFiveHoBase}{65.8}
\newcommand{\dialSeventyFiveHoFinal}{68.2}
\newcommand{\dialSeventyFiveStrongCalls}{29736}
\newcommand{\dialTwentyFiveAccEnd}{59.4}
\newcommand{\dialTwentyFiveAccStart}{51.5}
\newcommand{\dialTwentyFiveCalls}{9961}
\newcommand{\dialTwentyFiveGameEnd}{20.3}
\newcommand{\dialTwentyFiveGameStart}{23.8}
\newcommand{\dialTwentyFiveHoBase}{66.0}
\newcommand{\dialTwentyFiveHoFinal}{66.4}
\newcommand{\dialTwentyFiveStrongCalls}{9961}
\newcommand{\llamaCascadeAccEnd}{71.7}
\newcommand{\llamaCascadeAccStart}{62.0}
\newcommand{\llamaCascadeGameEnd}{12.7}
\newcommand{\llamaCascadeGameStart}{16.2}
\newcommand{\llamaCascadeHoBase}{53.6}
\newcommand{\llamaCascadeHoFinal}{63.0}
\newcommand{\llamaCascadeStrongCalls}{42134}
\newcommand{\llamaCheapAccEnd}{65.4}
\newcommand{\llamaCheapAccStart}{55.0}
\newcommand{\llamaCheapGameEnd}{18.2}
\newcommand{\llamaCheapGameStart}{21.8}
\newcommand{\llamaCheapHoBase}{53.2}
\newcommand{\llamaCheapHoFinal}{62.4}
\newcommand{\llamaCheapStrongCalls}{0}
\newcommand{\llamaSavings}{18}
\newcommand{\llamaStrongAccEnd}{73.5}
\newcommand{\llamaStrongAccStart}{62.7}
\newcommand{\llamaStrongGameEnd}{11.9}
\newcommand{\llamaStrongGameStart}{15.8}
\newcommand{\llamaStrongHoBase}{52.2}
\newcommand{\llamaStrongHoFinal}{63.6}
\newcommand{\llamaStrongStrongCalls}{51200}
\newcommand{\longCascadeAccEnd}{67.5}
\newcommand{\longCascadeAccStart}{58.0}
\newcommand{\longCascadeGameEnd}{12.8}
\newcommand{\longCascadeGameStart}{17.0}
\newcommand{\longCascadeHoBase}{66.2}
\newcommand{\longCascadeHoFinal}{69.2}
\newcommand{\longCascadeStrongCalls}{101504}
\newcommand{\longCheapAccEnd}{52.9}
\newcommand{\longCheapAccStart}{47.3}
\newcommand{\longCheapGameEnd}{29.1}
\newcommand{\longCheapGameStart}{28.7}
\newcommand{\longCheapHoBase}{65.8}
\newcommand{\longCheapHoFinal}{64.0}
\newcommand{\longCheapStrongCalls}{0}
\newcommand{\longSavings}{21}
\newcommand{\longStrongAccEnd}{66.6}
\newcommand{\longStrongAccStart}{59.5}
\newcommand{\longStrongBAccEnd}{66.4}
\newcommand{\longStrongBAccStart}{59.6}
\newcommand{\longStrongBGameEnd}{13.1}
\newcommand{\longStrongBGameStart}{17.6}
\newcommand{\longStrongBHoBase}{65.8}
\newcommand{\longStrongBHoFinal}{68.8}
\newcommand{\longStrongBStrongCalls}{128000}
\newcommand{\longStrongGameEnd}{12.8}
\newcommand{\longStrongGameStart}{16.5}
\newcommand{\longStrongHoBase}{65.8}
\newcommand{\longStrongHoFinal}{67.8}
\newcommand{\longStrongStrongCalls}{128000}
\newcommand{\mSevenBArmsComplete}{3}
\newcommand{\mSevenBCascadeAccEnd}{61.2}
\newcommand{\mSevenBCascadeAccStart}{55.0}
\newcommand{\mSevenBCascadeGameEnd}{17.9}
\newcommand{\mSevenBCascadeGameStart}{21.1}
\newcommand{\mSevenBCascadeHoBase}{65.8}
\newcommand{\mSevenBCascadeHoFinal}{68.2}
\newcommand{\mSevenBCascadeStrongCalls}{39503}
\newcommand{\mSevenBCheapAccEnd}{53.3}
\newcommand{\mSevenBCheapAccStart}{52.2}
\newcommand{\mSevenBCheapGameEnd}{26.1}
\newcommand{\mSevenBCheapGameStart}{23.4}
\newcommand{\mSevenBCheapHoBase}{66.0}
\newcommand{\mSevenBCheapHoFinal}{67.6}
\newcommand{\mSevenBCheapStrongCalls}{0}
\newcommand{\mSevenBSavings}{23}
\newcommand{\mSevenBStrongAccEnd}{62.4}
\newcommand{\mSevenBStrongAccStart}{55.1}
\newcommand{\mSevenBStrongGameEnd}{17.4}
\newcommand{\mSevenBStrongGameStart}{20.7}
\newcommand{\mSevenBStrongHoBase}{66.0}
\newcommand{\mSevenBStrongHoFinal}{67.8}
\newcommand{\mSevenBStrongStrongCalls}{51200}
\newcommand{\pilotThreeBCascadeAccEnd}{46.6}
\newcommand{\pilotThreeBCascadeAccStart}{45.4}
\newcommand{\pilotThreeBCascadeGameEnd}{13.9}
\newcommand{\pilotThreeBCascadeGameStart}{14.9}
\newcommand{\pilotThreeBCascadeHoBase}{59.4}
\newcommand{\pilotThreeBCascadeHoFinal}{61.4}
\newcommand{\pilotThreeBCascadeStrongCalls}{9605}
\newcommand{\pilotThreeBStrongAccEnd}{48.3}
\newcommand{\pilotThreeBStrongAccStart}{46.0}
\newcommand{\pilotThreeBStrongBaseAccEnd}{46.3}
\newcommand{\pilotThreeBStrongBaseAccStart}{46.2}
\newcommand{\pilotThreeBStrongBaseGameEnd}{14.2}
\newcommand{\pilotThreeBStrongBaseGameStart}{14.3}
\newcommand{\pilotThreeBStrongBaseHoBase}{58.8}
\newcommand{\pilotThreeBStrongBaseHoFinal}{57.8}
\newcommand{\pilotThreeBStrongBaseStrongCalls}{16000}
\newcommand{\pilotThreeBStrongGameEnd}{12.7}
\newcommand{\pilotThreeBStrongGameStart}{14.6}
\newcommand{\pilotThreeBStrongHoBase}{59.4}
\newcommand{\pilotThreeBStrongHoFinal}{60.0}
\newcommand{\pilotThreeBStrongStrongCalls}{16000}
\newcommand{\threeBSavings}{40}

% (auto-generated by ../verifier_economics_rlvr/scripts/make_rlvr_macros.py). [PENDING] = run
% not yet finished; NEVER hand-edit a measured number in here.
% ============================================================================================
\section{Cost-Repair in a Real RLVR Training Loop}\label{sec:rlvr}
The RNA/SAT testbed and the HumanEval+ study (\S\ref{sec:llm}) establish that a deterministic verifier is a
\emph{choice} with a measurable cost--robustness frontier, and that the cheapest-first cascade is provably
sound (Prop.~\ref{prop:cascade}). Those results are \emph{diagnostic}: they measure verifiers on fixed
candidate pools. The open question they raise---and the one a practitioner actually faces---is whether the
same object behaves as a \emph{priced, routable} reward when it sits \emph{inside} the RLVR optimization
loop, where the policy adapts to whatever the verifier accepts. This section closes that gap: we run GRPO
\citep{lambert2024tulu3} on mathematical reasoning with the verifier varied as the sole independent variable,
and ask what the choice of verifier buys and costs during training itself.

\paragraph{Setup.} The policy is Qwen2.5-Instruct \citep{qwen25_2024} at $3$B (pilot) and $7$B (main);
cross-family robustness is checked with Llama-3.1-8B-Instruct \citep{llama3_2024}. Training prompts are drawn
from NuminaMath-CoT \citep{numinamath2024}, restricted to items whose reference answer canonicalizes to a
plain numeric value so all verifiers are airtight; the held-out set is $500$ problems from the \texttt{math}
source, excluded from training entirely (MATH500-style). Because $60$--$90\%$ of GRPO groups on the raw pool
have zero reward variance (all-correct or all-wrong, hence zero gradient), the main runs use DAPO-style
difficulty filtering \citep{dapo2025}---keeping only prompts the base policy passes $1$--$7$ of $8$
times---which for the $7$B policy retains $7{,}113$ of $24{,}000$ prompts. All arms share model, data, seed,
optimizer, and step budget ($800$ steps, lr $2\!\times\!10^{-6}$); \textbf{the only difference is the reward
verifier.} We instrument three arms, exactly the rungs of \S\ref{sec:routing}:
\begin{itemize}\itemsep1pt
\item[] \textbf{cheap} --- any numeric token in the completion matches the gold value (over-accepts; gameable
by mentioning the right number while boxing another; near-free).
\item[] \textbf{strong} --- the boxed answer is value/symbolically equal to gold (objective-aware; paid on
every rollout).
\item[] \textbf{cascade} --- cheap gate, escalate to strong only on cheap-accepts. By the nesting
cheap $\supseteq$ strong, its reward equals strong's on every sample ($\epsilon{=}0$; Prop.~\ref{prop:cascade}),
while paying the strong cost only on cheap survivors.
\end{itemize}
Per rollout we log the cheap and strong verdicts, the gaming indicator (cheap-accept $\wedge$ strong-reject),
and verifier cost, from which we compute, per training step, true quality (strong-verified accuracy),
Youden's $J=\mathrm{TPR}-\mathrm{FPR}$ of the arm's accept decision against the strong gold, the gaming rate,
and cumulative strong-verifier calls. Held-out quality is the strong-verified accuracy of saved checkpoints
under greedy decoding.

\paragraph{Grounding the dial.} We take as given the phase-transition result of \citet{rateorfate2026}, who
show---via a replicator/bandit analysis of GRPO under a \emph{noisy} verifier---that the dynamics are governed
by Youden's index $J=\mathrm{TPR}-\mathrm{FPR}$: a sharp transition at $J{=}0$, with $J>0$ merely rescaling
\emph{convergence time} (``rate, not fate'') and $J<0$ amplifying incorrect reasoning modes to collapse.
Over-acceptance is the \emph{false-positive} direction of that verifier noise (it raises the FPR and thus
lowers $J$), so while $J>0$ it rescales convergence time rather than the final quality, with degradation only
once $J\le0$. This reframes the design target precisely---\emph{stay at
$J>0$ at the least verifier cost}---and predicts, rather than assumes, that a strictly-sound cheap cascade
should match the strong verifier's quality while a genuinely over-accepting cheap reward should converge
\emph{slower} (and, if pushed past $J\le0$, collapse). Our contribution is not that phase transition; it is
showing that the \emph{cascade / dial} realizes the favorable side of it \emph{at a measured fraction of the
verifier cost}.

\paragraph{Pilot (3B).} Across five $250$-step runs the loop behaves as the framework predicts on the axes it
can move: on-policy strong-accuracy rises ($+4$--$6$ points) while held-out accuracy is flat within noise---a
$250$-step, lr-$10^{-6}$ budget on $3$B is simply too short to move a $500$-item held-out set, which is why
the main runs lengthen the horizon and densify the signal. The load-bearing pilot result is on the
\emph{cost} axis and is model-agnostic: the cascade reproduces the strong arm's rewards exactly ($J=1.0$) while
issuing \textbf{\threeBSavings\%} fewer strong-verifier calls. This is the in-loop analogue of the diagnostic
cascade of \S\ref{sec:routing}, now measured against a moving policy.

\paragraph{Main results (7B).} Table~\ref{tab:rlvr7b} reports the three arms.
On the training distribution the arms separate exactly along the thesis: the strong and cascade arms improve
true (strong-verified) accuracy (\mSevenBStrongAccStart\%$\to$\mSevenBStrongAccEnd\% and
\mSevenBCascadeAccStart\%$\to$\mSevenBCascadeAccEnd\%) with gaming \emph{falling}, whereas the cheap arm's
inflated reward buys \emph{no} true-accuracy gain (\mSevenBCheapAccStart\%$\to$\mSevenBCheapAccEnd\%) and its
gaming rate \emph{rises} (\mSevenBCheapGameStart\%$\to$\mSevenBCheapGameEnd\%)---a clean in-loop
reward-hacking dissociation with the verifier as the sole manipulated variable. Crucially, the cascade tracks
the strong arm's true accuracy while issuing \textbf{\mSevenBSavings\%} fewer strong-verifier calls over
training (\mSevenBCascadeStrongCalls{} vs.\ \mSevenBStrongStrongCalls{})---buying the strong verifier's
decisions at a fraction of its price. We also observe a mechanism beyond mis-rewarding: because the cheap
reward saturates ($\sim\!0.76$ vs.\ $\sim\!0.44$ for strong/cascade), it drives most groups to zero reward
variance and thereby \emph{erases the very learning signal that difficulty filtering was constructed to
preserve}---a second, curriculum-level harm of over-acceptance.

\emph{Held-out generalization is flat within noise for \textbf{all} arms} at this budget (base$\to$final:
cheap \mSevenBCheapHoBase\%$\to$\mSevenBCheapHoFinal\%, strong \mSevenBStrongHoBase\%$\to$\mSevenBStrongHoFinal\%,
cascade \mSevenBCascadeHoBase\%$\to$\mSevenBCascadeHoFinal\%; $n{=}500$, $\pm\!\sim\!2$ points): even the strong
verifier does not move the general MATH held-out set in $800$ steps at lr $2\!\times\!10^{-6}$. We therefore
make \emph{no} held-out quality-repair claim---the arms' separation is on training dynamics (true-accuracy
gain, gaming) and on the \emph{cost} axis, not on held-out generalization. This is consistent with the $J>0$
regime of \citet{rateorfate2026}, where over-acceptance rescales convergence rather than the attained level;
extending the horizon or learning rate to where held-out moves for the strong arm---and confirming the cheap
arm then falls behind---is the natural next run.

\begin{table}[t]\centering\small
\setlength{\tabcolsep}{5pt}\renewcommand{\arraystretch}{1.15}
\caption{Main 7B runs (Qwen2.5-7B-Instruct, $800$ steps, difficulty-filtered NuminaMath; sole variable = the
reward verifier). On-policy = training-distribution strong-verified accuracy (first vs.\ last $100$ steps);
held-out = strict accuracy on $500$ MATH problems; strong-calls = cumulative strong-verifier invocations.
All values are computed from run logs by \texttt{make\_rlvr\_macros.py}, never entered by hand; cross-family
(Llama-3.1-8B) and second-seed rows are added the same way as those runs finish.}\label{tab:rlvr7b}
\begin{tabular}{lccccc}
\toprule
arm & on-policy acc & gaming & held-out (base$\to$final) & strong-calls & $J$\\
\midrule
cheap   & \mSevenBCheapAccStart\%$\to$\mSevenBCheapAccEnd\%     & \mSevenBCheapGameStart\%$\to$\mSevenBCheapGameEnd\%     & \mSevenBCheapHoBase\%$\to$\mSevenBCheapHoFinal\%     & \mSevenBCheapStrongCalls   & $\approx0.6$\\
strong  & \mSevenBStrongAccStart\%$\to$\mSevenBStrongAccEnd\%   & \mSevenBStrongGameStart\%$\to$\mSevenBStrongGameEnd\%   & \mSevenBStrongHoBase\%$\to$\mSevenBStrongHoFinal\%   & \mSevenBStrongStrongCalls  & $1.0$\\
cascade & \mSevenBCascadeAccStart\%$\to$\mSevenBCascadeAccEnd\% & \mSevenBCascadeGameStart\%$\to$\mSevenBCascadeGameEnd\% & \mSevenBCascadeHoBase\%$\to$\mSevenBCascadeHoFinal\% & \mSevenBCascadeStrongCalls & $1.0$\\
\bottomrule
\end{tabular}
\end{table}

\paragraph{Cost saving shrinks as the policy strengthens.} The cascade's saving is
$1-(\text{cheap-accept rate})$, and the cheap-accept rate rises as the policy learns to produce genuinely
right answers. Empirically the saving falls from $\sim\!40\%$ ($3$B) to \mSevenBSavings\% ($7$B): \emph{the
weaker the policy---exactly early training, where verifier cost dominates the budget---the more routing
saves.} We report this monotone dependence rather than a single headline multiplier, and make no
order-of-magnitude claim: in numeric-math the sound saving is a small constant, consistent with the RNA cost
finding (\S\ref{sec:cost}) and against the steep SAT case.

\paragraph{Honest scope.} That the cascade matches the strong arm's \emph{decisions} is guaranteed by
construction ($\epsilon{=}0$); it is not an empirical discovery and we do not present it as one. The empirical
contributions are three: (i) the \emph{magnitude and dynamics} of the cost saving in-loop, including its decay
with policy strength; (ii) that the cheap arm fails to convert inflated reward into true quality while gaming
rises, consistent with $J>0$ slow-convergence rather than immediate collapse \citep{rateorfate2026}; and
(iii) the curriculum-erasure mechanism. Cross-family (Llama-3.1-8B) and second-seed runs are in progress and
will populate the error bars in Table~\ref{tab:rlvr7b} and Figure~\ref{fig:rlvr} before submission.
